其上有 $1-\varphi$ 作用. 由蛇引理, 只需处理 $V_{\mathrm{ss}}$
与 $V/V_{\mathrm{ss}}$: 对前者, 利用由不动点组成的基, 命题~1.1 立即给出结论;
对后者, $\varphi$ 幂零, 因而 $1-\varphi$ 可逆.

\medskip
\noindent\emph{底本注.}
照录法文层保留了印本中 \emph{peut-être} 的连字符、单数措辞 \emph{coefficient tordus}、
\emph{Una application}、\emph{si elle additive} 中所漏的动词,
以及两处空格或标点冲突. 校订法文层和英文层规范了这些表面缺陷, 不改变数学内容.
公式 (1.0.7) 直接根据权威底本的高细节裁剪图读取: 其中的大运算符为交.

\section*{2. 凝聚上同调与平展上同调}

\subsubsection*{命题 2.0.}
设 $\bar X$ 为可分闭域 $k$ 上的固有概形, 设

\begin{equation}\tag{2.0.1}
\mathfrak F_p:H^i(\bar X,\mathcal O_{\bar X})
\longrightarrow H^i(\bar X,\mathcal O_{\bar X})
\end{equation}

为下式诱导的 $p$ 线性映射

\begin{equation}\tag{2.0.2}
\mathfrak F_p:\mathcal O_{\bar X}\longrightarrow\mathcal O_{\bar X},
\qquad \mathfrak F_p(f)=f^p.
\end{equation}

用

\[
H^i_{\mathrm{et}}(\bar X,\mathbb Z/p\mathbb Z)
\]

表示以常值层 $\mathbb Z/p\mathbb Z$ 为系数的平展上同调群. 则有典范同构

\begin{equation}\tag{2.0.3}
H^i_{\mathrm{et}}(\bar X,\mathbb Z/p\mathbb Z)
\simeq
H^i(\bar X,\mathcal O_{\bar X})^{1-\mathfrak F_p}.
\end{equation}

\medskip
\noindent\emph{证明.}
在 $\bar X_{\mathrm{et}}$ 上有层的正合列:

\begin{equation}\tag{2.0.4}
0\longrightarrow\mathbb Z/p\mathbb Z\longrightarrow\mathcal O_{\bar X}
\xrightarrow{\ 1-\mathfrak F_p\ }\mathcal O_{\bar X}\longrightarrow0,
\end{equation}

它给出上同调正合列

\begin{equation}\tag{2.0.5}
\cdots\longrightarrow H^i_{\mathrm{et}}(\bar X,\mathbb Z/p\mathbb Z)
\longrightarrow H^i_{\mathrm{et}}(\bar X,\mathcal O_{\bar X})
\xrightarrow{\ 1-\mathfrak F_p\ }
H^i_{\mathrm{et}}(\bar X,\mathcal O_{\bar X})\longrightarrow\cdots .
\end{equation}

有 $H^i_{\mathrm{et}}(\bar X,\mathcal O_{\bar X})\simeq H^i(\bar X,\mathcal O_{\bar X})$;
由于 $\bar X$ 固有, 这是有限维 $k$ 向量空间. 应用命题~1.2 得到正合列

\begin{equation}\tag{2.0.6}
0\longrightarrow H^i_{\mathrm{et}}(\bar X,\mathbb Z/p\mathbb Z)
\longrightarrow H^i(\bar X,\mathcal O_{\bar X})
\xrightarrow{\ 1-\mathfrak F_p\ }
H^i(\bar X,\mathcal O_{\bar X})\longrightarrow0,
\end{equation}

这就证明了断言. 现在命题~1.1.10 给出命题~2.0 的下述推论.

\subsubsection*{推论 2.1.}
对可分闭域 $k$ 上的固有概形 $\bar X$, 有典范同构

\begin{equation}\tag{2.1.1}
H^i(\bar X,\mathcal O_{\bar X})_{\mathrm{ss}}
\xleftarrow{\ \sim\ }
H^i_{\mathrm{et}}(\bar X,\mathbb Z/p\mathbb Z)
\otimes_{\mathbb F_p}k.
\end{equation}
